\documentclass[10pt,a4paper]{article}
\usepackage[top=1.4cm, bottom=1.4cm, left=1.6cm, right=1.6cm, headheight=14pt, footskip=18pt]{geometry}
\usepackage[utf8]{inputenc}
\usepackage{amsmath,amsfonts,amssymb}
\usepackage{graphicx}
\usepackage{tikz}
\usepackage{circuitikz}
\usetikzlibrary{calc,patterns}
\usepackage[most]{tcolorbox}
\usepackage{enumitem}
\usepackage{fancyhdr}
\usepackage{booktabs}
\usepackage{tabularx}
\usepackage{colortbl}
\usepackage{hyperref}

% --- Palet Warna Resmi UNIB ---
\definecolor{unibblue}{RGB}{0, 32, 96}
\definecolor{unibgold}{RGB}{197, 160, 89}
\definecolor{unibgreen}{RGB}{34, 112, 60}
\definecolor{boxbg}{RGB}{248, 250, 253}
\definecolor{linegray}{RGB}{210, 215, 225}

% --- Pengaturan Header & Footer ---
\pagestyle{fancy}
\fancyhf{}
\lhead{\footnotesize\textbf{\color{unibblue}Universitas Bengkulu} \textbar\ Program Studi S1 Teknik Elektro}
\rhead{\footnotesize\color{gray}Persamaan Diferensial (Genap 2026)}
\lfoot{\footnotesize\color{gray}Problem Set Tugas Terstruktur Berbasis OBE}
\cfoot{\footnotesize\color{gray}Hal. \thepage\ dari 2}
\rfoot{\footnotesize\color{unibblue}\textbf{Minggu 3: Aplikasi Transien Rekayasa \& Pem...}}
\renewcommand{\headrulewidth}{0.6pt}
\renewcommand{\footrulewidth}{0.4pt}

\begin{document}

% ==================== HALAMAN 1 ====================
\noindent\begin{minipage}[c]{0.68\textwidth}%
  {\large\textbf{\color{unibblue}PROBLEM SET (TUGAS MANDIRI TERSTRUKTUR)}}\\[1mm]
  {\normalsize\textbf{Mata Kuliah: Persamaan Diferensial (2 SKS)}}\\[0.5mm]
  {\small\textbf{Topik Minggu 3:} Aplikasi Transien Rekayasa \& Pemodelan Termal Transformator}%
\end{minipage}%
\hfill%
\begin{minipage}[c]{0.30\textwidth}%
  \begin{tcolorbox}[colback=unibblue!5,colframe=unibblue,boxrule=0.6pt,arc=2pt,left=4pt,right=4pt,top=2pt,bottom=2pt]
    \scriptsize
    \textbf{Nama:} \dotfill\\[1.2mm]
    \textbf{NPM:} \dotfill\\[1.2mm]
    \textbf{Kelas/Klp:} \dotfill\\[1.2mm]
    \textbf{Nilai Final:} \hfill \textbf{/ 100}
  \end{tcolorbox}%
\end{minipage}\par

\vspace{1.5mm}

% Kotak Sub-CPMK & Pedoman Tugas Terstruktur
\begin{tcolorbox}[colback=boxbg,colframe=unibgold!90!black,boxrule=0.6pt,arc=2pt,left=5pt,right=5pt,top=3pt,bottom=3pt,title=\textbf{\scriptsize Capaian Pembelajaran (Sub-CPMK 3) \& Pedoman Pengerjaan Tugas}]
  \scriptsize
  \textbf{Sub-CPMK 3:} Mahasiswa mampu menerapkan metodologi PDB orde 1 untuk memodelkan fenomena transien termal peralatan tenaga listrik (Hukum Pendinginan Newton \& IEEE Std C57.91) serta memvalidasinya secara numerik.\\
  \textbf{Ketentuan Tugas:} Kerjakan secara mandiri pada kertas folio/A4 bergaris atau laporan digital rapi. Lampirkan penurunan rumus analitis lengkap dan cetak visualisasi kode program Julia (Soal C6). Kumpulkan pada awal perkuliahan minggu berikutnya.
\end{tcolorbox}

\vspace{1.5mm}

\noindent{\textbf{\color{unibblue}\large BAGIAN I: FONDASI KONSEPTUAL \& PENURUNAN MATEMATIS}}\par
\vspace{1mm}

% Soal C1
\noindent\textbf{1. (C1 - Mengingat \textbar\ Bobot: 10 Poin)}\par\vspace{0.8mm}
{\footnotesize Tuliskan formulasi matematis Hukum Pendinginan Newton untuk temperatur benda $T(t)$ dalam medium sekitar $T_a$, serta sebutkan dimensi dan satuan SI untuk resistansi termal $R_{\text{th}}$ dan kapasitas kalor termal $C_{\text{th}}$.}\par

\vspace{2.5mm}

% Soal C2
\noindent\textbf{2. (C2 - Memahami \textbar\ Bobot: 15 Poin)}\par\vspace{0.8mm}
{\footnotesize Jelaskan analogi langsung antara parameter sistem kelistrikan (tegangan $V$, arus $I$, resistansi $R$, kapasitansi $C$) dengan parameter sistem termal (temperatur $T$, laju kalor $q$, resistansi termal $R_{\text{th}}$, kapasitas termal $C_{\text{th}}$).}\par

\vspace{2.5mm}

% Soal C3
\noindent\textbf{3. (C3 - Menerapkan \textbar\ Bobot: 20 Poin)}\par\vspace{0.8mm}
\noindent\begin{minipage}[t]{0.60\textwidth}%
  {\footnotesize Sebuah transformator distribusi gardu induk 20 kV memikul rugi daya tembaga $P_{\text{loss}} = 40\text{ kW}$. Sirkuit termal ekuivalen di samping memiliki kapasitas termal minyak $C_{\text{th}} = 2{,}0\text{ MJ/}^\circ\text{C}$ dan resistansi disipasi $R_{\text{th}} = 1{,}0 \times 10^{-3}\,^\circ\text{C/W}$. Jika suhu sekitar konstan $T_a = 32^\circ\text{C}$, tentukan fungsi kenaikan suhu minyak trafo $T(t)$ jika operasi dimulai dari keadaan dingin ($T(0) = T_a$).}%
\end{minipage}%
\hfill%
\begin{minipage}[t]{0.37\textwidth}%
  \centering
  \resizebox{0.95\linewidth}{!}{%
  \begin{circuitikz}[american, scale=0.68, transform shape]
    \draw (0,0) to[I, l=$P_{\text{loss}}$] (0,1.8)
          to[short] (1.6,1.8)
          to[R, l=$R_{\text{th}}$] (1.6,0)
          to[short] (0,0);
    \draw (1.6,1.8) to[short] (3.0,1.8)
          to[C, l=$C_{\text{th}}$, v=$\theta(t)$] (3.0,0)
          to[short] (1.6,0);
    \node at (3.0,2.1) [font=\tiny\bfseries, unibblue] {$T(t) - T_a$};
  \end{circuitikz}%
  }%
\end{minipage}\par

\vfill

% Kotak Formula Rujukan Teknis & Standar Industri
\begin{tcolorbox}[colback=unibblue!3,colframe=unibblue,colbacktitle=unibblue,coltitle=white,boxrule=0.6pt,arc=2pt,left=6pt,right=6pt,top=3pt,bottom=3pt,title=\textbf{\scriptsize FORMULA RUJUKAN TEKNIS \& STANDAR INDUSTRI TERKAIT}]
  \scriptsize
  \begin{itemize}[leftmargin=*, nosep]
  \item Hukum Pendinginan Newton: $\frac{dT}{dt} = -k (T - T_a)$ dengan $k = \frac{1}{R_{\text{th}} C_{\text{th}}}$
  \item Neraca Transien Termal Transformator: $C_{\text{th}} \frac{dT}{dt} + \frac{T - T_a}{R_{\text{th}}} = P_{\text{loss}}$
  \item Standar \textbf{IEEE Std C57.91}: Batas temperatur titik panas isolasi kertas-minyak: $T_{\text{kritis}} = 110^\circ\text{C}$
  \item Kenaikan Temperatur Tunak Minyak: $\Delta T_{\text{tunak}} = P_{\text{loss}} \cdot R_{\text{th}}$
\end{itemize}
\end{tcolorbox}

\newpage
% ==================== HALAMAN 2 ====================

\noindent{\textbf{\color{unibblue}\large BAGIAN II: ANALISIS SISTEM, EVALUASI STANDAR, \& KOMPUTASI JULIA}}\par
\vspace{1mm}

% Soal C4
\noindent\textbf{4. (C4 - Menganalisis \textbar\ Bobot: 20 Poin)}\par\vspace{0.8mm}
{\footnotesize Berdasarkan standar \textbf{IEEE Std C57.91}, suhu belitan transformator tidak boleh melampaui batas kritis isolasi kertas minyak $110^\circ\text{C}$. Jika transformator pada Soal 3 mengalami beban lebih mendadak sehingga rugi kalor meningkat menjadi $P_{\text{loss}} = 85\text{ kW}$, analisis berapa lama waktu operasi aman yang tersisa sebelum suhu mencapai batas kritis $110^\circ\text{C}$!}\par

\vspace{3.5mm}

% Soal C5
\noindent\textbf{5. (C5 - Mengevaluasi \textbar\ Bobot: 20 Poin)}\par\vspace{0.8mm}
{\footnotesize Sebuah sensor suhu RTD Pt100 diuji di laboratorium kalibrasi: saat dipindahkan dari bejana air mendidih ($100^\circ\text{C}$) ke penampung air es ($0^\circ\text{C}$), tercatat pembacaan suhu menjadi $36{,}8^\circ\text{C}$ pada detik ke-12. Evaluasi apakah sensor tersebut memiliki konstanta waktu termal $\tau = 12\text{ detik}$, dan tentukan pembacaan sensor pada detik ke-30!}\par

\vspace{3.5mm}

% Soal C6
\noindent\textbf{6. (C6 - Merancang \& Komputasi Julia \textbar\ Bobot: 15 Poin)}\par\vspace{0.8mm}
{\footnotesize Rancang skrip program ilmiah berbasis \textbf{Julia} (\texttt{Differential\-Equations.jl}) untuk memodelkan sistem pendinginan transformator dengan siklus pembebanan harian berfluktuasi: $P_{\text{loss}}(t) = 30 + 25\sin^2(\frac{\pi t}{43200})\text{ kW}$ selama 24 jam ($86400\text{ detik}$). Plot kurva suhu minyak $T(t)$ bersama batas aman IEEE C57.91.}\par

\vfill

% Tabel Rekapitulasi Skor OBE & Lembar Catatan Umpan Balik Dosen
\noindent\begin{minipage}[b]{0.62\textwidth}%
  \centering
  \scriptsize
  \begin{tabular}{|c|c|c|c|c|c|c|}
    \hline
    \rowcolor{unibblue}
    \textbf{\color{white}C1} & \textbf{\color{white}C2} & \textbf{\color{white}C3} & \textbf{\color{white}C4} & \textbf{\color{white}C5} & \textbf{\color{white}C6} & \textbf{\color{white}Total} \\
    \rowcolor{unibblue!10}
    10 Pts & 15 Pts & 20 Pts & 20 Pts & 20 Pts & 15 Pts & 100 Pts \\
    \hline
    & & & & & & \\[3mm]
    \hline
  \end{tabular}%
\end{minipage}%
\hfill%
\begin{minipage}[b]{0.35\textwidth}%
  \centering
  \scriptsize
  \textbf{Verifikasi Dosen / Asisten:}\\[6mm]
  (\dotfill)\\[1mm]
  \tiny Tanggal: \dotfill
\end{minipage}\par

\vspace{2mm}
\begin{tcolorbox}[colback=white,colframe=unibblue!40!black,colbacktitle=unibblue!10,coltitle=unibblue,boxrule=0.5pt,arc=2pt,left=5pt,right=5pt,top=3pt,bottom=3pt,title=\textbf{\tiny Catatan Evaluasi \& Umpan Balik Dosen Pengampu}]
  \tiny
  \textbf{Rubrik Pemeriksaan Pengerjaan Tugas Mandiri:}\par\vspace{0.8mm}
  $\square$ Penurunan matematis langkah-demi-langkah runtut (C1--C4) \hfill $\square$ Validasi analitis \& standar industri IEEE/IEC/SPLN akurat (C5)\\[0.8mm]
  $\square$ Skrip program Julia mandiri \& grafik visualisasi terlampir (C6) \hfill $\square$ Kerapian format laporan \& ketepatan waktu pengumpulan\\[2.0mm]
  \textbf{Catatan / Koreksi Khusus Dosen:}\\[1.6cm]
\end{tcolorbox}

\end{document}
